There is a long history of using self-play in multi-agent settings.
Early work explored self-play using genetic algorithms \citep{paredis1995,pollack1997,rosin1995,stanley2004}. \cite{sims1994a} and \cite{sims1994b} studied the emergent complexity in morphology and behavior of creatures that coevolved in a simulated 3D world. Open-ended evolution was further explored in the environments Polyworld \citep{yaeger1994} and Geb \citep{channon1998}, where agents compete and mate in a 2D world, and in Tierra \citep{ray1992} and Avida \citep{ofria2004}, where computer programs compete for computational resources. More recent work attempted to formulate necessary preconditions for open-ended evolution \citep{taylor2015, soros2014}. 
Co-adaptation between agents and environments can also give rise to emergent complexity \citep{florensa2017reverse, sukhbaatar2017intrinsic, wang2019}.
In the context of multi-agent RL, \cite{tesauro1995}, \cite{silver2016}, \cite{openai2018}, \cite{Jaderberg859} and \cite{vinyals2019} used self-play with deep RL techniques to achieve super-human performance in Backgammon, Go, Dota, Capture-the-Flag and Starcraft, respectively.
% \cite{bansal2017} trained agents to compete in a simulated 3D physics environment and found that the emergent behavior was more complex than the environment itself.
\cite{bansal2017} trained agents in a simulated 3D physics environment to compete in various games such as sumo wrestling and soccer goal shooting.
In \cite{liu2019emergent}, agents learn to manipulate a soccer ball in a 3D soccer environment and discover emergent behaviors such as ball passing and interception. In addition, communication has also been shown to emerge from multi-agent RL \citep{sukhbaatar2016, foerster2016, lowe2017, mordatch2017}.

Intrinsic motivation methods have been widely studied in the literature~\citep{chentanez2005intrinsically,singh2010intrinsically}.
One example is count-based exploration, where agents are incentivized to reach infrequently visited states by maintaining state visitation counts~\citep{strehl2008analysis,bellemare2016unifying,tang2017exploration} or density estimators~\citep{ostrovski2017count,burda2018exploration}.
Another paradigm are transition-based methods, in which agents are rewarded for high prediction error in a learned forward or inverse dynamics model \citep{schmidhuber1991possibility, stadie2015incentivizing, mohamed2015variational, houthooft2016vime, achiam2017surprise, pathak2017curiosity, burda2018large, haber2018learning}.  \citet{jacques2019} consider multi-agent scenarios and adopt causal influence as a motivation for coordination.
In our work, we utilize intrinsic motivation methods as an alternative exploration baseline to multi-agent autocurricula. Similar comparisons have also been made in \cite{haber2018learning} and \cite{leibo2019malthusian}.

Tool use is a hallmark of human and animal intelligence \citep{hunt1996manufacture, shumaker2011animal}; however, learning tool use in RL settings can be a hard exploration problem when rewards are unaligned.
For example, in \cite{forestier2017intrinsically, xie2019} a real-world robot learns to solve various tasks requiring tools. In \cite{bapst2019}, an agent solves construction tasks in a 2-D environment using both model-based and model-free methods. \cite{allen2019} uses a combination of human-designed priors and model-based policy optimization to solve a collection of physics-based puzzles requiring tool use.
However, in each of these works, agents were explicitly incentivized to interact with and use tools, whereas in our environment agents implicitly create this incentive through multi-agent competition.